%% kapa: exam=1 round=28 year=2017 period=2 form=A subject=accounting subject_ko=회계학 seq=15 number=55
%% source: Q-Net 감정평가사 기출문제 (cst00309) artlSeq=5199603 file=extracted/1008/5199603/2017년 제28회 감정평가사 1차시험 문제지/2교시 A형.pdf
%% conversion: kordoc 4.24.51 -> tools/gen1.py (auto); review=proofread-sonnet; flags=none
\begin{question}{55}
(주)감평은 20x1년 1월 1일 토지와 토지 위에 있는 건물A를 일괄하여 ￦40,000에 취득(토지와 건물A의 공정가치 비율은 4:1)하였다. 취득당시 건물A의 잔여 내용연수는 5년이고 잔존가치는 없으며 정액법으로 감가상각한다. 20x2년 1월 1일 더 이상 건물A를 사용할 수 없어 철거하고 새로운 건물B의 신축을 시작하였다. 건물A의 철거비용은 ￦1,500이며, 철거시 수거한 고철 등을 매각하여 ￦500을 수령하였다. 건물신축과 관련하여 20x2년에 ￦20,000의 건설비가 발생하였으며, 건물B(내용연수 10년, 잔존가치 ￦0, 정액법 감가상각)는 20x2년 10월 1일 완공 후 즉시 사용하였다. 20x1년 12월 31일 건물A의 장부금액과 20x2년 12월 31일 건물B의 장부금액은? (단, 감가상각은 월할계산한다.)
\begin{center}
\begin{tabular}{@{}lll@{\hspace{2.5em}}lll@{\hspace{2.5em}}lll@{}}
 & \uline{건물A} & \uline{건물B} &  & \uline{건물A} & \uline{건물B} &  & \uline{건물A} & \uline{건물B} \\
① & ￦6,400 & ￦19,500 & ② & ￦6,400 & ￦18,000 & ③ & ￦6,400 & ￦25,900 \\
④ & ￦8,000 & ￦19,500 & ⑤ & ￦8,000 & ￦26,900 &  &  &  \\
\end{tabular}
\end{center}
\end{question}
\answer{1}
